\section{Related Work}\label{sec:related}
% Draft 2026-10-17 (writer). Every statement about a cited work is checked against paper/notes/related_work.md,
% paper/notes/load_metrics_literature.md or docs/papers/dinos_claims.md. Self-citations in the third person
% (double-blind venues).

\paragraph{Autoscaling of cloud applications.}
Lorido-Botran et al.~\cite{lorido2014review} survey autoscaling techniques for elastic applications in the
cloud. In Kubernetes the standard mechanism is the HPA. It compares a metric such as CPU utilization with a
target, scales the replica count by their ratio, ignores changes within a tolerance, and holds back scale-down
decisions for a stabilization window~\cite{k8sdocs_hpa}. CPU of pods that are still starting is ignored for an
initialization period, so the controller sees a new pod's load only some time after the pod was
requested~\cite{k8sdocs_hpa}. AHPA, an autoscaler used in Alibaba Cloud's Kubernetes service, models the response
time of pods with M/M/1 and M/M/c queues and finds the M/M/c model more accurate when there are few
pods~\cite{zhou2023ahpa}. The AHPA description is a preprint. We use a threshold policy in the spirit of the HPA
as a baseline and a per-pod queueing model inside our simulator.

\paragraph{Deep reinforcement learning for autoscaling.}
DERP~\cite{bitsakos2018derp} applied deep reinforcement learning to elastic resource provisioning in the cloud.
DInos~\cite{bitsakos2025dinos} builds on it. Its agent reads a window of recent metrics (load, CPU, memory,
latency and the number of pods) through an LSTM~\cite{hochreiter1997lstm}, chooses every 15 seconds whether to
add a pod, remove one or do nothing, and is pre-trained in a simulator before it is fine-tuned in the target
environment, including a Kubernetes cluster serving Redis. Taylor and Stone~\cite{taylor2009transfer} survey
transfer in reinforcement learning and frame the question DInos asks: whether starting from a trained policy
beats training from scratch in the new task. Double Q-learning~\cite{vanhasselt2016ddqn} reduces the
overestimation of action values in deep Q-networks. Sofia and the retrained DInos baseline use its target.\todo{D8: confirm wording on how the DInos baseline is retrained (Double DQN target, fixes) before
Section~\ref{sec:agents} is written.} Sofia keeps the DInos state and action interface and changes the function
approximator. This lets us compare the two agents on the same inputs, the same actions and the same reward.

\paragraph{Performance and workload models for simulation.}
Queueing theory gives the standard link between load, capacity and response time~\cite{harcholbalter2013performance},
and Little's law~\cite{little1961proof} relates the number of requests in a system to their arrival rate and their
time in the system. For latency-critical services the tail matters more than the mean~\cite{dean2013tail}.
Measured tails of servers such as Memcached are worse than queueing models predict, because of interference from the
hardware, the operating system and the application~\cite{li2014tales}, and queueing and scheduling delays cause
violations of sub-millisecond targets at high utilization~\cite{leverich2014reconciling}. On the workload side,
Atikoglu et al.~\cite{atikoglu2012workload} analyse traces of a large Memcached deployment, Bodik et
al.~\cite{bodik2010spikes} characterize and generate workload spikes for stateful services, and Markov-modulated
Poisson processes~\cite{fischer1993mmpp} are a standard way to add bursts to a Poisson arrival stream. Production
traces such as the Azure Functions trace show daily and weekly patterns~\cite{shahrad2020serverless}. Scaling a
stateful store also has a cost of its own. Redis Cluster moves hash slots between nodes online, and keys are
locked briefly while they migrate~\cite{redisdocs_scaling,redisdocs_clusterspec}. Rocksteady shows that live
migration in an in-memory store can keep latency low under load~\cite{kulkarni2017rocksteady}. We use these works
to choose the structure of our simulator. Its parameters come from measurements on our own cluster, because we
found no published queueing model validated against Redis latency.\todo{Calibration pilot on the real cluster is
pending (D11); adjust this sentence if parameters stay provisional.}

\paragraph{Quantum reinforcement learning and hybrid models.}
Jerbi et al.~\cite{jerbi2021parametrized} train parametrized quantum policies with few qubits on benchmark
environments and prove a separation from classical learners on tasks constructed around the hardness of the
discrete logarithm. That separation concerns those constructed tasks and does not carry over to autoscaling.
Skolik et al.~\cite{skolik2022quantum} use VQCs as Q-functions in deep Q-learning. Their ablations show that
architectural and hyperparameter choices matter more than the parameter count, which is why we compare Sofia with
a classical control of the same size. Mari et al.~\cite{mari2020transfer} attach a variational circuit to a
pre-trained classical network, a form of transfer learning, and run it on quantum hardware for image and state
classification. Chen et al.~\cite{chen2020qlstm} place VQCs inside
the gates of an LSTM and report faster convergence in some of their tests. The work is a preprint, and the reported gains are specific
to those test cases. Bitsakos et al.~\cite{bitsakos2025quantuml} ran machine
learning algorithms, among them $k$-means, on a quantum cloud offering. None of these studies addresses
autoscaling, where the effect of an action is delayed, the reward trades cost against latency, and the policy must eventually run against a real
cluster within the control interval.\todo{Scholar: search for quantum RL applied to cloud or resource management
before any novelty statement is added here.}
